\subsection{Multivariate Polynomials}

In the previous sections, we worked with univariate polynomials $\mathbb{F}[X]$,
which are polynomials in a single variable. However, in a set of applications
such as \textit{Spartan}, we need to work with multivariate polynomials
$\mathbb{F}[X_1,\dots,X_v]$ where $v$ is the number of variables.

\begin{definition}
    We define a \textit{monomial} in $v$ variables as a product of the form:
    $X_1^{\alpha_1}\dots X_v^{\alpha_v}$ where $\alpha_1,\dots,\alpha_v$ are
    non-negative integers. The \textit{degree of a monomial} is defined as
    $\alpha_1+\dots+\alpha_v$. \textbf{Multivariate polynomial} from the space
    $\mathbb{F}[X_1,\dots,X_v]$ is defined as a finite linear combination of
    such monomials where coefficients are from $\mathbb{F}$. The \textbf{degree}
    of a multivariate polynomial is defined as the maximum degree of its
    monomials.
\end{definition}

\begin{example}
    For example, $f(X_1,X_2,X_3)=X_1^3 + 3X_1^2X_2^2 + X_3^2 + X_3$ is a linear
    combination of monomials $\{X_1^3,X_1^2X_2^2,X_3^2,X_3\}$, thus it is a
    multivariate polynomial in three variables $X_1$, $X_2$, and $X_3$. The
    maximum degree has a monomial $X_1^2X_2^2$ which is $4$, thus $\deg(f) = 4$.
\end{example}

However, in cryptography, we won't work with arbitrary multivariate
polynomials, compared to the univariate case. What we need is a 
\textit{multilinear polynomial}.
\begin{definition}
    A \textbf{multilinear polynomial} $f(X_1,\dots,X_v)$ is a multivariate
    polynomial which is linear in each variable, meaning that each variable
    appears with degree at most $1$. In other words, for each variable
    $X_i$, function $f(X_1,\dots,X_v)$ is a linear function for fixed 
    values of all other variables $X_1,\dots,X_{i-1},X_{i+1},\dots,X_v$:
    \begin{align*}
        f(X_1,\dots,X_v) &= \alpha_i(X_1,\dots,X_{i-1},X_{i+1},\dots,X_v) \cdot X_i \\
        &+ \beta_i(X_1,\dots,X_{i-1},X_{i+1},\dots,X_v)
    \end{align*}
\end{definition}

\begin{example}
    For example, $f(X_1,X_2,X_3)=X_1X_2 + 3X_1X_3 + X_2X_3$ is a multilinear 
    polynomial in three variables $X_1$, $X_2$, and $X_3$. Note that it is linear 
    in $X_1$. Indeed, it holds
    \begin{equation*}
        f(X_1,X_2,X_3) = \alpha_1(X_2,X_3)X_1 + \beta_1(X_2,X_3)
    \end{equation*}
    for $\alpha_1(X_2,X_3) = X_2 + 3X_3$ and $\beta_1(X_2,X_3) = X_2X_3$.

    Similarly, $f(X_1,\dots,X_v)=\prod_{i=1}^v X_i$ is also a multilinear
    polynomial in $v$ variables.
\end{example}

\subsection{Multilinear Extension}

Compared to univariate polynomials, multilinear polynomials can hold information
with a marge smaller degree. For example, if our interpolation domain is
$\{0,1,\dots,n-1\}$, then we need up to $n^{\text{th}}$ degree univariate
polynomial to interpolate the values over this domain. However, for multilinear
polynomials, we can encode the same domain with only $\log n$-variate
multilinear polynomial. Now, the details. 

\begin{definition}
    Suppose we are given the function $f: \{0,1\}^v \to \mathbb{F}$, mapping the
value from the $v$-dimensional hypercube to the field $\mathbb{F}$. We can
define the \textbf{extension} of $f$ as a $v$-variate polynomial $\widetilde{f}
\in \mathbb{F}[X_1,\dots,X_v]$ which agrees with $f$ on the points of the
hypercube, meaning that for each $\boldsymbol{b} \in \{0,1\}^v$ it holds that
$f(\boldsymbol{b}) = \widetilde{f}(\boldsymbol{b})$.
\end{definition}

\begin{example}
    A more intuitive way to think about the extension is that given $2^v$ points
    of the hypercube, we construct the ``interpolation'' function
    $\widetilde{f}$ which agrees with the given values over the hypercube. For
    example, if $f(0,0) = 1$, $f(0,1) = 2$, $f(1,0) = 2$, and $f(1,1) = 3$, then
    the extension $\widetilde{f}$ can be defined as
    $\widetilde{f}(X_1,X_2)=X_1^2+X_2^2+1$.
\end{example}

There are possibly very large number of possible extensions for a given 
function $f$. What we are interested in is the multilinear extension, which
turns out to be unique.

\begin{theorem}\label{theorem:multilinear-extension}
    Any function over the $v$-dimensional hypercube $f: \{0,1\}^v \to
    \mathbb{F}$ has a unique $v$-variate multilinear extension $\widetilde{f}
    \in \mathbb{F}[X_1,\dots,X_v]$. It is defined using the \textit{Lagrange
    interpolation of multilinear polynomials} formula as follows:
    \begin{equation*}
        \widetilde{f}(\mathbf{X}) = \sum_{\boldsymbol{b} \in \{0,1\}^v} f(\boldsymbol{b}) \cdot \mathsf{eq}(\mathbf{X};\boldsymbol{b}),
    \end{equation*}
    where the set $\{\mathsf{eq}(\mathbf{X}; \boldsymbol{b})\}_{\boldsymbol{b} \in
    \{0,1\}^v}$ is referred to as \textit{the set of multilinear Lagrange basis
    polynomials} over the set $\{0,1\}^v$. Each
    $\mathsf{eq}(\mathbf{X};\boldsymbol{b})$ is defined as:
    \begin{equation*}
        \mathsf{eq}(\mathbf{X};\boldsymbol{b}) \triangleq \prod_{i=1}^v \{X_ib_i + (1-X_i)(1-b_i) \}.
    \end{equation*}
\end{theorem}

\textbf{Proof.} First, let us show that the provided $\widetilde{f}$ is indeed a
multilinear polynomial. Notice that each $\mathsf{eq}(\mathbf{X};\boldsymbol{b})$ is
a multilinear polynomial, since it is a product of linear polynomials in each
variable $X_i$. Thus, $\widetilde{f}$ is a linear combination of multilinear
polynomials, hence it is also a multilinear polynomial. 

Why is it an extension? Notice that the formula for $\mathsf{eq}(\mathbf{X};
\boldsymbol{b})$ has the following property: $\mathsf{eq}(\mathbf{X}; \boldsymbol{b})=1$
if $\mathbf{X}=\boldsymbol{b}$, and $0$ otherwise. Indeed, if $X_i \neq b_i$, then
the term $X_ib_i + (1-X_i)(1-b_i)$ is equal to $0$, and thus the whole product
is equal to $0$. If $X_i = b_i$ for each $i$, then each term is equal to $1$ and
thus the product is $1$. Therefore, $\widetilde{f}(\boldsymbol{b}')$ for each
$\boldsymbol{b}' \in \{0,1\}^v$ is equal to:
\begin{equation*}
    \widetilde{f}(\boldsymbol{b}') = \sum_{\boldsymbol{b} \in \{0,1\}^v} f(\boldsymbol{b}) \cdot \mathsf{eq}(\boldsymbol{b}';\boldsymbol{b}) = f(\boldsymbol{b}') \cdot \underbrace{\mathsf{eq}(\boldsymbol{b}';\boldsymbol{b}')}_{=1} + \sum_{\boldsymbol{b} \neq \boldsymbol{b}'} f(\boldsymbol{b}) \cdot \underbrace{\mathsf{eq}(\boldsymbol{b}';\boldsymbol{b})}_{=0} = f(\boldsymbol{b}').
\end{equation*}

\begin{example}\label{example:mle}
    Suppose function $f: \{0,1\}^2 \to \mathbb{F}_p$ for $p=11$ is given by
    $f(0,0)=3$, $f(0,1)=4$, $f(1,0)=1$, and $f(1,1)=2$. We first build the
    multilinear Lagrange basis polynomials:
    \vspace{-7.5px}
    \begin{gather*}
        \mathsf{eq}(X_1,X_2;(0,0)) = (1-X_1)(1-X_2), \quad
        \mathsf{eq}(X_1,X_2;(0,1)) = (1-X_1)X_2, \\
        \mathsf{eq}(X_1,X_2;(1,0)) = X_1(1-X_2), \quad
        \mathsf{eq}(X_1,X_2;(1,1)) = X_1X_2.
    \end{gather*}
    The multilinear extension $\widetilde{f}(X_1,X_2)$ is thus computed as:
    \vspace{-7.5px}
    \begin{equation*}
        \widetilde{f}(X_1,X_2) = 3(1-X_1)(1-X_2) + 4(1-X_1)X_2 + X_1(1-X_2) + 2X_1X_2 = 3 - 2X_1 + X_2.
    \end{equation*}
\end{example}

Now, the question is, how fast can we compute the multilinear extension
$\widetilde{f}$ given $2^v$ values of $f$? Consider the following lemma.

\begin{lemma}\label{lemma:multilinear-extension-computation}
    Fix some positive integer $v$ and let $n=2^v$. Given an input
    $f(\boldsymbol{b})$ for all $\boldsymbol{b} \in \{0,1\}^v$ and a vector
    $\boldsymbol{r}=(r_1,\dots,r_v) \in \mathbb{F}^{\log n}$, one can compute
    $\widetilde{f}(\mathbf{r})$ in $\mathcal{O}(n)$ time and space.
\end{lemma}

\subsection{The Sum-Check Protocol}

\subsubsection{Protocol Description}

Suppose we are given the $v$-variate polynomial (possibly non-multilinear) $f:
\{0,1\}^v \to \mathbb{F}$ over a finite field $\mathbb{F}$. The main goal
of the Sum-Check protocol \cite{sumcheck} is to convince the verifier $\mathcal{V}$ that
\begin{equation*}
    \sum_{b_1 \in \{0,1\}}\sum_{b_2 \in \{0,1\}} \dots \sum_{b_v \in \{0,1\}} f(b_1,\dots,b_v) = H
\end{equation*}
for the given value $H \in \mathbb{F}$. In other words, we can convince that the
sum of all the values of $f$ over the hypercube $\{0,1\}^v$ (which we further
write as $\sum_{\boldsymbol{b} \in \{0,1\}^v}f(\boldsymbol{b})$ for short) is
equal to the given value $H$. Such check might be useless at first glance, but
as it turns out many protocols can be reduced to the Sum-Check protocol,
similarly how any NP statement can be encoded as a high-degree univariate
polynomial check (see the lecture on QAP).

\textbf{Why can't the verifier just compute the sum?} Typically, $v$ is a fairly
large number. According to the \Cref{lemma:multilinear-extension-computation},
the verifier can compute the multilinear extension $\widetilde{f}$ in
$\mathcal{O}(2^v)$ time and space, which is infeasible for large $v$. The
Sum-Check protocol allows the verifier to check the sum in up to
$\mathcal{O}(v^2)$ time and space, which is much more adequate for large $v$.

\textbf{How does it work?} First, the prover $\mathcal{P}$ sends the value $C_1
\in \mathbb{F}$, which he claims to be the value of the sum (that is, $H$).
The protocol proceeds in $v$ rounds. For each round $j$, define the 
following univariate polynomial in the variable $X_j$:
\begin{equation*}
    f_j(X_j) = \sum_{(b_{j+1},\dots,b_v) \in \{0,1\}^{v-j}} f(r_1,\dots,r_{j-1},X_j,b_{j+1},\dots,b_v),
\end{equation*}
where values $r_1,\dots,r_{j-1} \in \mathbb{F}$ are fixed values of the
variables $X_1,\dots,X_{j-1}$ (which are randomnesses selected during previous
rounds).

Consider the first round, when $j=1$. In such case, according to the definition
of $f_j(X_j)$, the prover $\mathcal{P}$ computes the univariate polynomial
$f_1(X_1) = \sum_{(b_2,\dots,b_v) \in \{0,1\}^{v-1}} f(X_1,b_2,\dots,b_v)$ and 
sends the \textit{claimed} polynomial $s_1(X_1)$ as the first round message.
How can the verifier $\mathcal{V}$ be sure that $s_1(X_1)$ is indeed the
univariate polynomial $f_1(X_1)$? Since $f_1(0) + f_1(1) = H$, the verifier
can check that $s_1(0) + s_1(1) = C_1$. 

However, this is not enough: there are many univariate polynomials that satisfy
the condition $s_1(0) + s_1(1) = C_1$. For that reason, we are going to apply
the Schwartz-Zippel lemma, which states that as long as $|\mathbb{F}| \gg \deg
f_1$, it is safe to check the equality $s_1(r_1) = f_1(r_1)$ at a random point
$r_1 \leftarrowS \mathbb{F}$. In such case, the soundness of such check is
$1-\deg f_1/|\mathbb{F}|$. 

However, can the verifier evaulate both $s_1(r_1)$ and $f_1(r_1)$ effectively to
verify the equality? Good news is that $s_1(r_1)$ can be evaluated efficiently: in
fact, in $\mathcal{O}(\deg s_1)$ where typically degree is small (in case of
multilinear polynomials, the degree is at most $1$). But what about $f_1(r_1)$?
Here is the bad news: $f_1(r_1)$ is a sum of $2^{v-1}$ terms, and thus it cannot
be computed efficiently. Another good news is that we do not need to! The idea
of the sumcheck protocol is to reduce the computation of $f_1(r_1)$ to the
computation of $f_2$, then $f_3$ etc. until the computation is trivial. 

For concrete example, consider the second round. Now, the prover computes the
polynomial $f_2(X_2) = \sum_{(b_3,\dots,b_v) \in \{0,1\}^{v-2}}
f(r_1,X_2,b_3,\dots,b_v)$ and sends the claimed polynomial $s_2(X_2)$. The
verifier now checks that $s_2(0) + s_2(1) = s_1(r_1)$ and that $s_2(r_2) =
f_2(r_2)$ at random point $r_2 \leftarrowS \mathbb{F}$. Since computing
$f_2(r_2)$ consists in adding up $2^{v-2}$ terms, it is still infeasible to
compute it directly, but we reduced the problem $4 \times$ already.

\textbf{Recursive definition.} For the $j$-th round, the prover computes
$f_j(X_j)$ and sends the claimed polynomial $s_j(X_j)$. The verifier checks
whether $s_j(0) + s_j(1) = s_{j-1}(r_{j-1})$. Additionally, the verifier rejects
if $\deg s_j$ is too large (for example, if $\deg s_j > \deg_j f$ where $\deg_j
f$ is a degree of the polynomial in variable $X_j$). During the last round, the
prover sends the claimed value $s_v(r_v)$ and the verifier checks whether
$s_v(r_v) = f(r_1,\dots,r_v)$\footnote{Here we assume that the verifier has the
oracle access to the function $f$ and can compute it at any randomly selected
point.}. If the check succeeds, then the verifier accepts the claim that
\begin{equation*}
    \sum_{\boldsymbol{b} \in \{0,1\}^v} f(\boldsymbol{b}) = H.
\end{equation*}

One might ask: how secure is the Sum-Check protocol? Consider the following 
lemma.
\begin{lemma}\label{lemma:sumcheck-protocol}
    Let $f \in \mathbb{F}[X_1,\dots,X_v]$ be a multivariate polynomial of degree
    at most $d$ in each variable, defined over the finite field $\mathbb{F}$.
    For any given $H \in \mathbb{F}$, let $\mathcal{L}$ be the language of 
    all polynomials $f$ (given as an oracle) such that
    \begin{equation*}
        H = \sum_{b_1 \in \{0,1\}}\sum_{b_2 \in \{0,1\}} \dots \sum_{b_v \in \{0,1\}} f(b_1,\dots,b_v).
    \end{equation*}
    The sumcheck protocol is an IOP for $\mathcal{L}$ with the completness error
    $\delta_C=0$ and the soundness error $\delta_S \leq vd/|\mathbb{F}|$.
\end{lemma}

Finally, consider the efficiency of the Sum-Check protocol.
\begin{lemma}
    Suppose $f$ is given according to the representation in Lemma \ref{lemma:sumcheck-protocol}. Then, the 
    following holds:
    \begin{itemize}
        \item Communication consists of $\mathcal{O}(dv)$ field elements.
        \item The verifier $\mathcal{V}$ runs in $\mathcal{O}(vd)+T$ where $T$
        is the cost of the oracle access to $f$.
        \item The prover $\mathcal{P}$ runs in $\mathcal{O}(2^vT)$.
    \end{itemize}
\end{lemma}

\subsubsection{Sum-Check Protocol Implementation}

Let us implement this algorithm in SageMath! First, define the multivariate
polynomial ring and the function $f: \{0,1\}^v \to \mathbb{F}$. As an example,
we will use $v=10$ while $f$ will be sampled randomly from
$\mathbb{F}_p[X_1,\dots,X_v]$. As a prime field, we use $p=2^{31} - 1$. 
Here is the code that sets everything up:
\begin{minted}{python}
# Defining the finite field GF(p) where p is a large prime
p = (1<<31) - 1
Fp = GF(p)

# Defining the multivariate polynomial ring over the finite field GF(p)
v = 10 # Degree of the polynomial
variable_names = [f'x{i}' for i in range(v)]
R = PolynomialRing(Fp, names=variable_names)
variables = R.gens()
\end{minted}

Now, let us define $f$ and find the corresponding value $H =
\sum_{\boldsymbol{b} \in \{0,1\}^v} f(\boldsymbol{b})$. For that, we
simply take the random element of the polynomial ring \textsf{R} and 
manually add up all $2^v$ values over the boolean hypercube $\{0,1\}^v$.
\begin{minted}{python}
def boolean_hypercube_sum(f: PolynomialRing) -> Fp:
    """
    Computes the sum of the polynomial f over the boolean hypercube {0,1}^v.
    The boolean hypercube is represented by evaluating the polynomial at all
    combinations of 0 and 1 for each variable.
    """

    total_sum = Fp(0)
    for i in range(1<<v):
        # Convert i to a binary representation of length v
        binary_representation = [(i >> j) & 1 for j in range(v)]
        # Evaluate the polynomial at this point
        point_value = f(*binary_representation)
        total_sum += point_value
    return total_sum

f = R.random_element(degree=v) # Random polynomial
H = boolean_hypercube_sum(f)
\end{minted}

The SageMath generated the following $f$ and $H$:
\begin{align*}
    f(X_1,\dots,X_{10}) &= 226921892X_2^3X_3X_4^2X_5X_6X_8^2 + 274566947X_0X_1^2X_2^2X_4X_7^2X_8X_9 \\ 
    &+ 850030718X_1^2X_4X_6^3X_7^2X_8X_9 + 560601732X_0X_3X_5X_6X_7X_8X_9^4 - 390725845X_0^2X_2^3X_3X_9^3 \\
    H &= \sum_{\boldsymbol{b} \in \{0,1\}^{10}} f(\boldsymbol{b}) = 1053620759.
\end{align*}

Now, we are going to implement the prover algorithm. For that, we need to define
the univariate polynomial $f_j(X_j)$ for each round $j$ and compute the claimed
polynomial $s_j(X_j)$ based on previously sampled randomness values
$r_1,\dots,r_{j-1}$. Since our protocol is a \textit{public-coin} protocol, we
can apply the \textit{Fiat-Shamir transform} to make the randomness values
$r_1,\dots,r_v$ deterministically computable by the prover (thus making the
protocol non-interactive). For that, we will compute the challenge $r_j$ as
$\mathcal{H}(s_1, s_2, \dots, s_{j-1})$ where $\mathcal{H}: \{0,1\}^* \to
\mathbb{F}_p$ is a cryptographic hash function (SHA256 in our case). That 
said, we define $\mathcal{H}$ as follows:
\begin{minted}{python}
def fiat_shamir_sample(input_string: str) -> Fp:
    """
    Samples a random value from the finite field GF(p) using the Fiat-Shamir heuristic.
    The input string is hashed to produce a random value.
    """

    fs_hash = hashlib.sha256(input_string.encode('utf-8')).hexdigest()
    return Fp(int(fs_hash, 16) % p)
\end{minted}

Finally, the verifier logic is even more simple: the verifier checks that (a)
randomnesses $\{r_j\}_{j \in \{1,\dots,10\}}$ are sampled correctly, (b) the
claimed polynomials $s_j(X_j)$ satisfy the conditions $s_j(0) + s_j(1) =
s_{j-1}(r_{j-1})$ for each $j \in \{1,\dots,10\}$ (instead of $s_0$, we simply
use the claimed value $H$), and (c) the final check $s_v(r_v) =
f(r_1,\dots,r_v)$ holds. Additionally, we should check the degrees of each
$s_j$, but we omit this check for simplicity. Here is the implementation of the
prover and verifier:
\begin{minted}{python}
class SumCheckProtocol:
    """
    Prover and Verifier for the Sum Check protocol.
    """

    def __init__(
        self, 
        polynomial: PolynomialRing, 
        claimed_sum: Fp
    ) -> None:
        """
        Initializes the SumCheck protocol with a polynomial and the claimed sum.

        Args:
            polynomial (PolynomialRing): The polynomial to be checked.
            claimed_sum (Fp): The sum claimed by the prover over the boolean hypercube.
        """

        self.f = polynomial
        self.H = claimed_sum


    @staticmethod
    def fiat_shamir_from_polynomials(polynomials: list[PolynomialRing]) -> Fp:
        """
        Samples a random value from the finite field GF(p) using the Fiat-Shamir heuristic
        based on the provided polynomials.

        Args:
            polynomials (list[PolynomialRing]): List of polynomials to be used for sampling.

        Returns:
            Fp: A random value from the finite field GF(p).
        """
        
        fs_input = ''.join(str(poly) for poly in polynomials)
        return fiat_shamir_sample(fs_input)


    def prove(self) -> dict:
        """
        Prover generates the transcript of the polynomial and the claimed sum.

        Returns:
            dict: A transcript containing the polynomial, claimed sum, random values, and polynomials.
        """

        # Initialize the transcript representing the 
        # interaction between the prover and verifier.
        transcript = {
            'polynomial': self.f,
            'claimed_sum': self.H,
            'random_values': [],
            'polynomials': []
        }

        for j in range(v):
            # The protocol consists of v rounds.
            # Finding polynomial s_j that represents the sum of f over the boolean hypercube
            Rj = PolynomialRing(Fp, 'x')
            x = Rj.gen()
            s_j = Rj.zero()

            for i in range(2**(v-j-1)):
                # Convert i to a binary representation of length v-j-1
                binary_representation = [(i >> k) & 1 for k in range(v-j-1)]
                # Evaluate the polynomial at this point
                point_value = self.f(*transcript['random_values'][:j], variables[j], *binary_representation)
                univariate_poly = point_value.subs({variables[j]: x})
                s_j += univariate_poly

            # Append the polynomial s_j to the transcript
            transcript['polynomials'].append(s_j)

            # Sample the random value using Fiat-Shamir heuristic
            random_value = SumCheckProtocol.fiat_shamir_from_polynomials(
                polynomials=transcript['polynomials']
            )
            transcript['random_values'].append(random_value)

        return transcript


    def verify(self, transcript: dict) -> bool:
        """
        Verifier checks the validity of the transcript.

        Args:
            transcript (dict): The transcript generated by the prover.

        Returns:
            bool: True if the claimed sum is valid, False otherwise.
        """
        
        # Decipher the transcript
        H = transcript['claimed_sum']
        f = transcript['polynomial']
        random_values = transcript['random_values']
        polynomials = transcript['polynomials']

        # Assert that random_values is formed correctly using Fiat-Shamir heuristic
        if len(random_values) != v:
            print("Invalid number of random values in the transcript.")
            return False
        for i in range(v):
            if random_values[i] != self.fiat_shamir_from_polynomials(polynomials[:i+1]):
                print(f"Random value at index {i} does not match the Fiat-Shamir sample.")
                return False

        # Assert that each s_r(0) + s_r(1) matches the claimed sum H
        for j in range(v):
            # During the first round, simply check that s_1(0) + s_1(1) = H
            if j == 0:
                s_1 = polynomials[0]
                if s_1(0) + s_1(1) != H:
                    print(f"""First round sum {s_1(0) + s_1(1)} 
                            does not match the claimed sum {H}.""")
                    return False
                continue
            
            # For subsequent rounds, check that the sum of the polynomials matches the claimed sum
            s, s_previous = polynomials[j], polynomials[j-1]
            if s(0) + s(1) != s_previous(random_values[j-1]):
                print(f"""Round {j} sum {s(0) + s(1)} does not match the 
                        previous round's output {s_previous(random_values[j-1])}.""")
                return False

        # Final round checks whether f(r) = s_v(r_v)
        last_polynomial = polynomials[-1]
        last_random = random_values[-1]

        if last_polynomial(last_random) != f(*random_values):
            print(f"Final check failed: {last_polynomial(last_random)} != {f(*random_values)}")
            return False
        
        return True

protocol = SumCheckProtocol(f, H)
transcript = protocol.prove()
print(f'Prover has the transcript: {transcript}')
print(f'Verifier checks the proof: {protocol.verify(transcript)}')
\end{minted}

% Finally, we can run the protocol as follows:
% \begin{minted}{python}
% protocol = SumCheckProtocol(f, H)
% transcript = protocol.prove()
% print(f'Prover has the transcript: {transcript}')
% print(f'Verifier checks the proof: {protocol.verify(transcript)}')
% \end{minted}

The example transcript is the following:
\begin{align*}
    s_1(X) &= 763349684X^2 + 238898491X + 25686292, \quad r_1 = 493136960 \\
    s_2(X) &= 1430156880X^2 + 361631142, \quad r_2 = 2831006 \\
    s_3(X) &= 658473518X^3 + 208437747X^2 + 1648222287, \quad r_3 = 321757611 \\
    s_4(X) &= 111234379X + 1501686780, \quad r_4 = 1835658 \\
    s_5(X) &= 2135686754X^2 + 1501686780X + 1698421434, \quad r_5 = 1970078146 \\
    s_6(X) &= 53232575X + 259099570, \quad r_6 = 1616339175 \\
    s_7(X) &= 669263002X^3 + 514137765X + 1942401931, \quad r_7 = 887816643 \\
    s_8(X) &= 1524551309X^2 + 1551518026X + 55160592, \quad r_8 = 421872749 \\
    s_9(X) &= 1766229428X^2 + 826393776X + 1291949229, \quad r_9 = 1169032581 \\
    s_{10}(X) &= 2030644729X^4 + 1291949229X^3 + 560585988X + 1720313399, \\
    r_{10} &= 1461328437
\end{align*}

\subsection{Sumcheck Applications}

As of now, the Sum-Check protocol seems too abstract and not very useful: why do
we even need to sum some random multivariate polynomial over the boolean
hypercube? In this section, we provide some (rather theoretic) applications of
the Sum-Check protocol, but as we go further, it will become more clear why it
is so important.

\subsubsection{The \#SAT Problem}

Let $\mathtt{C}: \{0,1\}^{\ell} \to \{0,1\}$ be any boolean formula of size $S =
\mathcal{O}(\mathsf{poly}(\ell))$. In the \#SAT problem, the goal is to compute
the number of satisfying (boolean) assignments of $\mathtt{C}$, that is, find
the value $H = \sum_{\boldsymbol{b} \in
\{0,1\}^{\ell}}\mathtt{C}(\boldsymbol{b})$. Such problem is believed to be very
difficult with the fastest known algorithm to still run in the exponential time:
that is, no much better than brute-force the formula in time
$\mathcal{O}(2^{\ell}S)$. Even determining whether $H > 0$ is widely believed to
be NP-hard. But suppose we have a prover $\mathcal{P}$ who does know $H$ and
wants to prove any verifier $\mathcal{V}$ that $H$ was computed correctly in the
polynomial time.

Now, similarly to how it is done in R1CS, we \textit{arithmetize} the circuit
$\mathtt{C}$ to be computable over the finite field $\mathbb{F}$ (let us call it
$\widetilde{\mathtt{C}}: \mathbb{F}^{\ell} \to \mathbb{F}$). Here how it goes:
\begin{itemize}
    \item Instead of the gate $x \wedge y$, we use the multiplication $x \cdot y$.
    \item Instead of the gate $x \vee y$, we use the addition $x + y - x \cdot y$.
    \item Instead of the gate $\neg x$, we use the subtraction $1 - x$.
    \item Instead of the gate $x \oplus y$, we use the addition $x + y - 2xy$.
\end{itemize}

As an example, suppose we have the circuit:
\begin{equation*}
    \mathtt{C}(x_1,x_2,x_3,x_4) = (\neg x_1 \wedge x_2) \wedge (x_3 \vee x_4)
\end{equation*}

The extension $\widetilde{\mathtt{C}}$ of the circuit $\mathtt{C}$ is
\begin{equation*}
    \widetilde{\mathtt{C}}(X_1,X_2,X_3,X_4) = (1 - X_1)X_2(X_3 + X_4 - X_3X_4).
\end{equation*}

It is fairly easy to see that $\widetilde{\mathtt{C}}(\boldsymbol{b}) =
\mathtt{C}(\boldsymbol{b})$ for all $\boldsymbol{b} \in \{0,1\}^{\ell}$ with
$\ell=4$. Then, we can apply the sum-check protocol over
$\widetilde{\mathtt{C}}$ to prove that $H$ was computed correctly. In such case,
the prover runs in $\mathcal{O}(S^22^{\ell})$ time, while the verifier in time
$\mathcal{O}(S)$. 

\subsubsection{Matrix Multiplication Verification (MatMul Check-Sum Protocol)}

Now, here is the application which gets much more interesting. Consider two
matrices $A, B \in \mathbb{F}^{n \times n}$ and the goal is to verify that the
product $C = A \cdot B$ is computed correctly. The naive way is to make verifier
compute the product, taking $\mathcal{O}(n^3)$ time, and then compare the result
with the claimed value $C$. However, it is possible to verify the correctness of
the product in $\mathcal{O}(n^2)$ time and space using the \textit{Freivelds'}
protocol by the simple observation that we can sample a random challenge
$\boldsymbol{\alpha} \leftarrowS \mathbb{F}^n$ and verify that
$A(B\boldsymbol{\alpha}) = C\boldsymbol{\alpha}$ (idea is pretty similar
to the Schwartz-Zippel lemma where we check the polynomial equation at a 
random point). However, is there any better way to do that? Sumcheck 
protocol can help to keep the same assymptotics, but do not reveal 
the whole matrices $A,B$ to the verifier.

\textbf{Protocol.} Here is where \textit{multilinear extension} comes into play.
Instead of perceiving matrices $A,B,C$ as $n^2$ field elements, we perceive 
them as functions $f_A, f_B, f_C: \{0,1\}^{\log n} \times \{0,1\}^{\log n} \to \mathbb{F}$,
mapping two indices of the matrix to the corresponding value. This way, for instance:
\begin{equation*}
    f_A(\boldsymbol{i}, \boldsymbol{j}) = A_{i,j}, \quad \text{where} \quad \boldsymbol{i}=(i_1,\dots,i_{\log n}), \; \boldsymbol{j}=(j_1,\dots,j_{\log n}).
\end{equation*}

Denote by $\widetilde{f}_A, \widetilde{f}_B, \widetilde{f}_C: \mathbb{F}^{\log n} \times \mathbb{F}^{\log n} \to \mathbb{F}$ the
multilinear extensions of the functions $f_A, f_B, f_C$. Now, how to 
reduce the seemingly difficult check $C = AB$ into the Sum-Check
protocol? Consider the following lemma.
\begin{lemma}
    $\widetilde{f}_C(\mathbf{x},\mathbf{y}) = \sum_{\boldsymbol{b} \in \{0,1\}^{\log n}} \widetilde{f}_A(\mathbf{x}, \boldsymbol{b})\widetilde{f}_B(\boldsymbol{b}, \mathbf{y})$.
\end{lemma}

\textbf{Proof.} Obviously, both left and right-hand sides are multilinear
polynomials in $\mathbf{x}$ and $\mathbf{y}$. Since multilinear extension of $C$
is unique, it suffices to check that the equality holds for all boolean
assignments $\boldsymbol{i}, \boldsymbol{j} \in \{0,1\}^{\log n}$. Indeed, we
have $\widetilde{f}_C(\boldsymbol{i},\boldsymbol{j}) = \sum_{\boldsymbol{b} \in
\{0,1\}^{\log n}} \widetilde{f}_A(\boldsymbol{i},
\boldsymbol{b})\widetilde{f}_B(\boldsymbol{b}, \boldsymbol{j})$. But notice that
this check literally checks whether $C_{i,j} = \sum_{b=1}^n A_{i,b}B_{b,j}$
which is exactly the definition of matrix multiplication!

Now, the interactive proof is immediate: sample random $\boldsymbol{r}_1$,
$\boldsymbol{r}_2 \leftarrowS \mathbb{F}^{\log n}$, and apply the sum-check
on $g(\mathbf{z}) := \widetilde{f}_A(\boldsymbol{r}_1,
\mathbf{z})\widetilde{f}_B(\mathbf{z}, \boldsymbol{r}_2)$ to prove that it 
equals $\widetilde{f}_C(\boldsymbol{r}_1, \boldsymbol{r}_2)$.

It can be shown that both the prover's and verifier's time is
$\mathcal{O}(n^2)$, while the proof size is $\mathcal{O}(\log n)$ (compared to
$\mathcal{O}(n^2)$ for the naive approach).

\subsubsection{MatMul Sum-Check Protocol Implementation}

Now, let us implement the matrix multiplication verification protocol in
SageMath! For simplicity, we assume that the matrix has a size $n=2^v$ so that
MLE of matrices are of degree $2v$. Also, we will work over the small prime
field $\mathbb{F}_{11}$ this time to make outputs somewhat more readable. 
So here we go:
\begin{minted}{python}
# Defining the finite field GF(p) where p is a large prime
p = 11
Fp = GF(p)

# Assume for simplicity that matrices are of size 2^v x 2^v
v = 2 # Matrix of size 4x4
n = 1<<v

# Defining how MLE are computed
variable_names = [f'x{i}' for i in range(2*v)]
R = PolynomialRing(Fp, names=variable_names)
variables = R.gens()
\end{minted}

Now, for concreteness, we initialize the matrices $A,B$ and $C=AB$.
\begin{minted}{python}
# Generate two random matrices A and B over the finite field GF(p)
A = Matrix(Fp, n, n, [Fp.random_element() for _ in range(n*n)])
B = Matrix(Fp, n, n, [Fp.random_element() for _ in range(n*n)])

# Find the product matrix C=A*B
C = A * B
print(f'Matrix A:\n{A}')
print(f'Matrix B:\n{B}')
print(f'Matrix C:\n{C}')
\end{minted}

Here, we obtained the following matrices:
\begin{equation*}
    A = \begin{bmatrix}
    1 & 1 & 3 & 10 \\
    5 & 6 & 4 & 5 \\
    1 & 8 & 3 & 4 \\
    6 & 2 & 2 & 7
    \end{bmatrix}, \quad B = \begin{bmatrix}
    2 & 7 & 9 & 5 \\
    5 & 10 & 0 & 10 \\
    5 & 0 & 2 & 3 \\
    2 & 6 & 6 & 6
    \end{bmatrix}, \quad C = \begin{bmatrix}
    9 & 0 & 9 & 7 \\
    4 & 4 & 6 & 6 \\
    10 & 1 & 6 & 8 \\
    2 & 5 & 1 & 10
    \end{bmatrix}
\end{equation*}

Now, we need to build the multilinear extensions $\widetilde{f}_A,
\widetilde{f}_B, \widetilde{f}_C: \mathbb{F}^{2 \times 2} \to \mathbb{F}$. To
achieve that, we first write down the function to compute the MLE given the set
of values $\{(\boldsymbol{b}, f(\boldsymbol{b}))\}_{\boldsymbol{b} \in
\{0,1\}^{\ell}}$ given interpolation formula from
Theorem~\ref{theorem:multilinear-extension}.
\begin{minted}{python}
def mle_from_hypercube(hypercube: list) -> R:
    """
    Computes the Multivariate Linear Extension (MLE) of a hypercube.
    The hypercube is represented as a list of tuples, where each tuple
    contains the coordinates of a point in the hypercube.

    Args:
        hypercube (list): A list of tuples representing points in the hypercube.
                           Each tuple should have length equal to the dimension dim.
    """

    # Create a polynomial for each point in the hypercube
    mle = R.zero()
    for point, value in hypercube:
        eq_poly = R(1)
        for i, bit in enumerate(point):
            eq_poly *= bit*variables[i] + (1-bit)*(1-variables[i])

        mle += eq_poly * value
    
    return mle
\end{minted}

We can verify the correctness of this function by running the Example in the
first section. Indeed:
\begin{minted}{python}
print('Example MLE:', mle_from_hypercube([
    ((0, 0), Fp(3)),
    ((0, 1), Fp(4)),
    ((1, 0), Fp(1)),
    ((1, 1), Fp(2))
])) # Output: -2*x0 + x1 + 3
\end{minted}

At this point, we can finally derive the MLEs. We simply iterate through all
indices $(i,j)$, bit-decompose them to get bit-vectors
$\boldsymbol{i}=(i_1,i_2)$ and $\boldsymbol{j}=(j_1,j_2)$, and then interpolate
the MLEs: for matrix $A$, for instance, we have $\{(i_1,i_2,j_1,j_2),
A_{i,j}\}_{i,j \in [n]^2}$.
\begin{minted}{python}
def mle_from_matrix(matrix: Matrix) -> R:
    """
    Computes the Multivariate Linear Extension (MLE) of a matrix.
    The MLE is a polynomial that represents the matrix entries as variables.
    """
    
    assert matrix.nrows() == matrix.ncols(), "Matrix must be square."
    assert matrix.nrows() == n, "Matrix size must be 1<<v"

    # Range over all indices, bit-decompose them, and build the mle
    hypercube = []
    for i in range(n):
        for j in range(n):
            # Convert i and j to binary representation of length v
            point = [(i >> k) & 1 for k in range(v)] + [(j >> k) & 1 for k in range(v)]
            hypercube.append((tuple(point), matrix[i, j]))

    return mle_from_hypercube(hypercube)
\end{minted}

Now we find the MLEs of matrices $A$, $B$, and $C$:
\begin{minted}{python}
A_mle = mle_from_matrix(A)
B_mle = mle_from_matrix(B)
C_mle = mle_from_matrix(C)
print(f'MLE of A:\n{A_mle}')
print(f'MLE of B:\n{B_mle}')
print(f'MLE of C:\n{C_mle}')
\end{minted}

The results are the following:
\begin{align*}
\widetilde{f}_A(X_1,X_2,Y_1,Y_2) &= -X_1X_2Y_1 - 3X_1X_2Y_2 + 4X_1Y_1Y_2 - 2X_2Y_1Y_2 \\
& \quad + X_1X_2 + X_1Y_1 - 4X_2Y_1 - 3X_1Y_2 - 4Y_1Y_2 \\
& \quad + 4X_1 + 2Y_2 + 1 \\
\widetilde{f}_B(X_1,X_2,Y_1,Y_2) &= -2X_1X_2Y_1Y_2 - 2X_1X_2Y_1 - 3X_1X_2Y_2 + 3X_1Y_1Y_2 \\
& \quad + 4X_2Y_1Y_2 + 5X_1X_2 + X_2Y_1 - X_1Y_2 + X_2Y_2 \\
& \quad + 2Y_1Y_2 + 3X_1 + 3X_2 + 5Y_1 - 4Y_2 + 2 \\
\widetilde{f}_C(X_1,X_2,Y_1,Y_2) &= 2X_1X_2Y_1Y_2 + 3X_1X_2Y_1 + X_1X_2Y_2 + 4X_1Y_1Y_2 \\
& \quad + 4X_2Y_1Y_2 - 3X_1X_2 - 2X_1Y_1 + 2X_1Y_2 - 4X_2Y_2 \\
& \quad - 4Y_1Y_2 - 5X_1 + X_2 + 2Y_1 - 2
\end{align*}

One can easily check that they indeed encode the matrices $A$, $B$, and $C$ as
expected. For example, it is trivial to check the upper left element in all 
three cases:
\begin{equation*}
    \widetilde{f}_A(0,0,0,0) = 1, \quad \widetilde{f}_B(0,0,0,0) = 2, \quad \widetilde{f}_C(0,0,0,0) = 9.
\end{equation*}

Now, we can implement the Sum-Check protocol to verify that $C = AB$ holds. We
first sample randomnesses $\boldsymbol{r}_1, \boldsymbol{r}_2 \leftarrowS
\mathbb{F}_{11}^2$:
\begin{minted}{python}
# Now, we are going to apply the sumcheck protocol. First, 
# sample two random vectors r1 and r2 of size v = log(n)
r1 = [Fp.random_element() for _ in range(v)]
r2 = [Fp.random_element() for _ in range(v)]
\end{minted}

Then, we define the function $g(\mathbf{z}) = \widetilde{f}_A(\boldsymbol{r}_1,
\mathbf{z})\widetilde{f}_B(\mathbf{z}, \boldsymbol{r}_2)$ and apply the
Sum-Check protocol to prove that it equals $\widetilde{f}_C(\boldsymbol{r}_1,
\boldsymbol{r}_2)$. For that, we first define the function $g$:
\begin{minted}{python}
claimed_sum = C_mle(*(r1 + r2))
g = A_mle.subs({
    variables[i]: r1[i] for i in range(v)
}).subs({
    variables[v+i]: variables[i] for i in range(v)
}) * B_mle.subs({
    variables[i+v]: r2[i] for i in range(v)
})
\end{minted}

Turnes out that the claimed sum is $H=6$ while the polynomial is:
\begin{equation*}
    g(X_1,X_2) = 5X_1^2X_2^2 - 5X_1^2X_2 + 3X_1X_2^2 + 3X_1^2 + 4X_1X_2 - 3X_2^2 + 4X_1 - 3X_2 + 2
\end{equation*}

So the only thing left is to run the already implemented Sum-Check protocol.
Here how it goes:
\begin{minted}{python}
# Defining a smaller polynomial ring for the protocol (since g 
# is polynomial in v variables instead of 2*v)
Q = PolynomialRing(Fp, names=variable_names)
variables = Q.gens()
g = Q(g)

# Running the sum-check protocol
protocol = SumCheckProtocol(Fp, Q, g, claimed_sum, degree=v)
transcript = protocol.prove()
print_transcript(transcript)
verification_result = protocol.verify(transcript)
print(f'Verification result: {verification_result}')
\end{minted}

Our interaction consisted of only two polynomials (since $v=2$):
\begin{align*}
    s_1(X) &= 6X^2 + 4X + 9, &&& r_1 &= 4 \\
    s_2(X) &= X^2 + 10X, &&& r_2 &= 2
\end{align*}

The final verification result is \textsf{True}, meaning prover has successfully
proved that $C = AB$ holds. 
